\documentclass{beamer}


\usepackage[T2A]{fontenc}
\usepackage[utf8]{inputenc}
\usepackage[english,russian]{babel}
\pdfmapfile{+sansmathaccent.map}


\mode<presentation>
{
	\usetheme{Warsaw} % or try Darmstadt, Madrid, Warsaw, Rochester, CambridgeUS, ...
	\usecolortheme{seahorse} % or try seahorse, beaver, crane, wolverine, ...
	\usefonttheme{serif}  % or try serif, structurebold, ...
	\setbeamertemplate{navigation symbols}{}
	\setbeamertemplate{caption}[numbered]
} 


%%%%%%%%%%%%%%%%%%%%%%%%%%%%
% itemize settings


%%%%%%%%%%%%%%%%%%%%%%%%%%%%
% itemize settings

\definecolor{myhotpink}{RGB}{255, 80, 200}
\definecolor{mywarmpink}{RGB}{255, 60, 160}
\definecolor{mylightpink}{RGB}{255, 80, 200}
\definecolor{mypink}{RGB}{255, 30, 80}
\definecolor{mydarkpink}{RGB}{155, 25, 60}

\definecolor{mypaleblue}{RGB}{240, 240, 255}
\definecolor{mylightblue}{RGB}{120, 150, 255}
\definecolor{myblue}{RGB}{90, 90, 255}
\definecolor{mygblue}{RGB}{70, 110, 240}
\definecolor{mycourseblue}{RGB}{120, 140, 255}
\definecolor{mycourselightblue}{RGB}{245, 250, 255}
\definecolor{mydarkblue}{RGB}{0, 0, 180}
\definecolor{myblackblue}{RGB}{40, 40, 120}

\definecolor{mygreen}{RGB}{0, 200, 0}
\definecolor{mygreen2}{RGB}{245, 255, 230}
\definecolor{mydarkgreen}{RGB}{0, 120, 0}
\definecolor{mydarkgreen}{RGB}{0, 120, 0}


\definecolor{mygray}{gray}{0.8}
\definecolor{mydarkgray}{RGB}{80, 80, 160}
\definecolor{mylightgray}{RGB}{160, 160, 160}

\definecolor{mydarkred}{RGB}{160, 30, 30}
\definecolor{mylightred}{RGB}{255, 150, 150}
\definecolor{myred}{RGB}{200, 110, 110}
\definecolor{myblackred}{RGB}{120, 40, 40}

\definecolor{mypink}{RGB}{255, 30, 80}
\definecolor{myhotpink}{RGB}{255, 80, 200}
\definecolor{mywarmpink}{RGB}{255, 60, 160}
\definecolor{mylightpink}{RGB}{255, 80, 200}
\definecolor{mydarkpink}{RGB}{155, 25, 60}
\definecolor{mywhitepink}{RGB}{255, 240, 240}

\definecolor{mydarkcolor}{RGB}{60, 25, 155}
\definecolor{mylightcolor}{RGB}{130, 180, 250}



\setbeamertemplate{itemize items}[default]

\setbeamertemplate{itemize item}{\color{myblackblue}$\blacksquare$}
\setbeamertemplate{itemize subitem}{\color{mygblue}$\blacktriangleright$}
\setbeamertemplate{itemize subsubitem}{\color{mygray}$\blacksquare$}

\setbeamercolor{palette quaternary}{fg=white,bg=mycourseblue}
\setbeamercolor{titlelike}{parent=palette quaternary}

\setbeamercolor{palette quaternary2}{fg=black,bg=mycourselightblue}
\setbeamercolor{frametitle}{parent=palette quaternary2}

\setbeamerfont{frametitle}{size=\Large,series=\scshape}
\setbeamerfont{framesubtitle}{size=\normalsize,series=\upshape}





%%%%%%%%%%%%%%%%%%%%%%%%%%%%
% block settings

\setbeamercolor{block title}{bg=red!30,fg=black}

\setbeamercolor*{block title example}{bg=mygreen!40!white,fg=black}

\setbeamercolor*{block body example}{fg= black, bg= mygreen2}


%%%%%%%%%%%%%%%%%%%%%%%%%%%%
% URL settings
\hypersetup{
	colorlinks=true,
	linkcolor=blue,
	filecolor=blue,      
	urlcolor=blue,
}

%%%%%%%%%%%%%%%%%%%%%%%%%%

\renewcommand{\familydefault}{\rmdefault}


\newcommand{\mycolor}[1] {
\setbeamercolor{normal text}{fg={#1}} % Recolors standard body text
\setbeamercolor{structure}{fg={#1}}   % Recolors bullets and structure items
\usebeamercolor[fg]{normal text}      % Applies the normal text color immediately
}


\usepackage{amsmath}
\usepackage{mathtools}

\usepackage{cancel}
\usepackage{centernot}

\usepackage{changepage}
\usepackage{subcaption}

\usepackage{qrcode}

\newcommand{\bo}[1] {\mathbf{#1}}
\newcommand{\R}{\mathbb{R}} 
\newcommand{\T}{^\top}     



\newcommand{\mydate}{Fall 2026}

\newcommand{\mygit}{\textcolor{blue}{\href{https://github.com/SergeiSa/Advanced-Control-2026}{github.com/SergeiSa/Advanced-Control-2026}}}

\newcommand{\myqr}{ \textcolor{black}{\qrcode[height=1.5in]{https://github.com/SergeiSa/Advanced-Control-2026}}
}




\newcommand{\bref}[2] {\textcolor{blue}{\href{#1}{#2}}}

%%%%%%%%%%%%%%%%%%%%%%%%%%%%
% code settings

\usepackage{listings}
\usepackage{color}
% \definecolor{mygreen}{rgb}{0,0.6,0}
% \definecolor{mygray}{rgb}{0.5,0.5,0.5}
\definecolor{mymauve}{rgb}{0.58,0,0.82}
\lstset{ 
	backgroundcolor=\color{white},   % choose the background color; you must add \usepackage{color} or \usepackage{xcolor}; should come as last argument
	basicstyle=\footnotesize,        % the size of the fonts that are used for the code
	breakatwhitespace=false,         % sets if automatic breaks should only happen at whitespace
	breaklines=true,                 % sets automatic line breaking
	captionpos=b,                    % sets the caption-position to bottom
	commentstyle=\color{mygreen},    % comment style
	deletekeywords={...},            % if you want to delete keywords from the given language
	escapeinside={\%*}{*)},          % if you want to add LaTeX within your code
	extendedchars=true,              % lets you use non-ASCII characters; for 8-bits encodings only, does not work with UTF-8
	firstnumber=0000,                % start line enumeration with line 0000
	frame=single,	                   % adds a frame around the code
	keepspaces=true,                 % keeps spaces in text, useful for keeping indentation of code (possibly needs columns=flexible)
	keywordstyle=\color{blue},       % keyword style
	language=Octave,                 % the language of the code
	morekeywords={*,...},            % if you want to add more keywords to the set
	numbers=left,                    % where to put the line-numbers; possible values are (none, left, right)
	numbersep=5pt,                   % how far the line-numbers are from the code
	numberstyle=\tiny\color{mygray}, % the style that is used for the line-numbers
	rulecolor=\color{black},         % if not set, the frame-color may be changed on line-breaks within not-black text (e.g. comments (green here))
	showspaces=false,                % show spaces everywhere adding particular underscores; it overrides 'showstringspaces'
	showstringspaces=false,          % underline spaces within strings only
	showtabs=false,                  % show tabs within strings adding particular underscores
	stepnumber=2,                    % the step between two line-numbers. If it's 1, each line will be numbered
	stringstyle=\color{mymauve},     % string literal style
	tabsize=2,	                   % sets default tabsize to 2 spaces
	title=\lstname                   % show the filename of files included with \lstinputlisting; also try caption instead of title
}


%%%%%%%%%%%%%%%%%%%%%%%%%%%%
% URL settings
\hypersetup{
	colorlinks=false,
	linkcolor=blue,
	filecolor=blue,      
	urlcolor=blue,
}

%%%%%%%%%%%%%%%%%%%%%%%%%%

%%%%%%%%%%%%%%%%%%%%%%%%%%%%
% tikz settings

\usepackage{tikz}
\tikzset{every picture/.style={line width=0.75pt}}


\newcommand{\myqrframe}{
	\begin{frame}
		\begin{center}
			\centerline{Слайды доступны на Github:}
			\bigskip
			\centerline{\mygit}
			\bigskip
			\myqr
		\end{center}
	\end{frame}
}


\subtitle{Современные методы управления роботами}
\author{Сергей Савин}
\date{2026}

\title{Линейный-квадратичный регулятор, Уравнение Беллмана}
\begin{document}
	\maketitle
	
	
	\begin{frame}{Содержание}
		\begin{itemize}
			\item Оптимальность и функция ценности
			\item Уравнение Беллмана
			\item Конечно-горизонтный дискретный LQR
			\begin{itemize}
				\item Оптимальный регулятор
				\item Рекурсия Риккати
			\end{itemize}
			\item Бесконечно-горизонтный дискретный LQR
			\begin{itemize}
				\item Алгебраическое уравнение Риккати
				\item Численное решение уравнения Риккати
			\end{itemize}
		\end{itemize}
	\end{frame}
	
	
	
	%------------------------------------------------
	\begin{frame}{Оптимальность}
		\begin{flushleft}
			
			Рассмотрим дискретную динамическую систему:
			\[
			x_{k+1}=f(x_k,u_k).
			\]
			
			Закон управления задаётся в виде
			\[
			u_k=\pi(x_k).
			\]
			
			Для конечного горизонта в $N$ шагов пусть стоимость равна
			\[
			J=\phi(x_N)+\sum_{k=0}^{N-1}\ell(x_k,u_k).
			\]
			
			\medskip
			
			Теперь мы можем поставить следующий вопрос:
			
			\medskip
			
			\emph{Каков оптимальный закон управления $\pi^*$ и
				соответствующая траектория, минимизирующие стоимость $J$?}
			
		\end{flushleft}
	\end{frame}
	%------------------------------------------------
	
	\begin{frame}{Функция ценности}
		\begin{flushleft}
			
			Определим функцию ценности $V_k(x)$ как минимальную
			стоимость, достижимую при старте из $x_k=x$ с
			учётом оставшихся $N-k$ шагов:
			\[
			\begin{aligned}
				V_k(x)=\min_{u_k,\ldots,u_{N-1}}\quad
				&\phi(x_N)+\sum_{i=k}^{N-1}\ell(x_i,u_i)\\
				\text{subject to}\quad
				&x_{i+1}=f(x_i,u_i),\\
				&x_k=x.
			\end{aligned}
			\]
			
			\medskip
			
			Функция ценности представляет собой \textbf{оптимальную
				стоимость до конца}: наименьшую возможную стоимость при
			фиксированном числе оставшихся шагов.
			
		\end{flushleft}
	\end{frame}
	%------------------------------------------------
	
	\begin{frame}{Уравнение Беллмана}
		\begin{flushleft}
			
			Уравнение Беллмана описывает, как функция ценности
			на одном шаге связана с функцией ценности на
			следующем шаге:
			\[
			\boxed{
				V_k(x)=\min_u
				\left[
				\ell(x,u)+V_{k+1}\bigl(f(x,u)\bigr)
				\right].
			}
			\]
			
			\medskip
			
			Оптимальная стоимость до конца состоит из двух частей:
			\begin{itemize}
				\item Мгновенная стоимость $\ell(x,u)$.
				\item Оптимальная стоимость до конца из следующего состояния.
			\end{itemize}
			
			\medskip
			
			На терминальном шаге времени не остаётся, поэтому
			функция ценности равна терминальной стоимости:
			\[
			V_N(x)=\phi(x).
			\]
			
		\end{flushleft}
	\end{frame}
	%------------------------------------------------
	
	
	
	
	%------------------------------------------------
	\begin{frame}{Линейный случай}
		\begin{flushleft}
			
			Рассмотрим линейную динамику:
			\[
			x_{k+1}=Ax_k+Bu_k.
			\]
			
			Пока предположим, что закон управления линейный:
			\[
			u_k=-Kx_k.
			\]
			
			Динамика замкнутой системы имеет вид
			\[
			x_{k+1}=(A-BK)x_k.
			\]
			
			\medskip
			
			Рассмотрим квадратичную функцию стоимости:
			\[
			J=x_N^\top Q_Nx_N+
			\sum_{i=0}^{N-1}
			\left(x_i^\top Qx_i+u_i^\top Ru_i\right),
			\]
			где $Q_N\succeq0$, $Q\succeq0$ и $R\succ0$.
			
		\end{flushleft}
	\end{frame}
	%------------------------------------------------
	
	\begin{frame}{Линейный случай: функция ценности}
		\begin{flushleft}
			{\small
				Сначала рассмотрим фиксированный линейный регулятор
				$u_k=-Kx_k$.
				
				При этом регуляторе будущее состояние зависит от
				текущего состояния следующим образом:
				\[
				x_i=(A-BK)^{i-k}x_k.
				\]
				
				Кроме того,
				\[
				x_i^\top Qx_i+u_i^\top Ru_i
				=x_i^\top(Q+K^\top RK)x_i.
				\]
				
				\medskip
				
				Следовательно, стоимость до конца при этом регуляторе
				является квадратичной: $V_k(x)=x^\top S_kx$, где
				\[
				\begin{aligned}
					S_k={}&
					\bigl((A-BK)^{N-k}\bigr)^\top
					Q_N(A-BK)^{N-k}\\
					&+\sum_{i=k}^{N-1}
					\bigl((A-BK)^{i-k}\bigr)^\top
					(Q+K^\top RK)
					(A-BK)^{i-k}.
				\end{aligned}
				\]
				
				\medskip
				
				Та же квадратичная форма будет описывать оптимальную
				функцию ценности, но с матрицами $S_k$, определяемыми
				уравнением Беллмана.}
			
		\end{flushleft}
	\end{frame}
	%------------------------------------------------
	
	
	
	%------------------------------------------------
	\begin{frame}{Линейный случай: уравнение Беллмана}
		\begin{flushleft}
			
			Уравнение Беллмана имеет вид
			\[
			V_k(x)=\min_u
			\left[
			\ell(x,u)+V_{k+1}(Ax+Bu)
			\right].
			\]
			
			Для квадратичной стоимости предположим, что
			\[
			V_k(x)=x^\top S_kx.
			\]
			
			Тогда уравнение Беллмана принимает вид
			\[
			x_k^\top S_kx_k
			=\min_{u_k}\left[
			x_k^\top Qx_k+u_k^\top Ru_k
			+x_{k+1}^\top S_{k+1}x_{k+1}
			\right].
			\]
			
			\medskip
			
			Поскольку $x_{k+1}=Ax_k+Bu_k$, можно записать,
			используя $x=x_k$ и $u=u_k$:
			\[
			\boxed{
				x^\top S_kx=\min_u
				\left[
				x^\top Qx+u^\top Ru
				+(Ax+Bu)^\top S_{k+1}(Ax+Bu)
				\right]
			}
			\]
			
		\end{flushleft}
	\end{frame}
	%------------------------------------------------
	
	\begin{frame}{Оптимальный регулятор}
		\begin{flushleft}
			
			Чтобы найти оптимальное управление, продифференцируем
			выражение внутри минимизации по $u$
			и приравняем производную к нулю:
			\[
			\begin{aligned}
				&2Ru+2B^\top S_{k+1}Bu +2B^\top S_{k+1}Ax=0.
			\end{aligned}
			\]
			
			Эквивалентно,
			\[
			(R+B^\top S_{k+1}B)u
			=-B^\top S_{k+1}Ax.
			\]
			
			\medskip
			
			Разрешая относительно $u$, получаем оптимальный закон управления:
			\[
			\boxed{
				u=-K_kx
			}
			\]
			где
			\[
			\boxed{
				K_k=(R+B^\top S_{k+1}B)^{-1}
				B^\top S_{k+1}A
			}
			\]
			
			\medskip
			
			Чтобы вычислить оптимальный коэффициент усиления на каждом шаге,
			нам нужна матрица $S_{k+1}$.
			
		\end{flushleft}
	\end{frame}
	%------------------------------------------------
	
	\begin{frame}{Дискретная рекурсия Риккати}
		\begin{flushleft}
			
			Подстановка оптимального управления в уравнение Беллмана
			даёт
			\[
			\begin{aligned}
				x^\top S_kx
				={}&x^\top Qx+x^\top A^\top S_{k+1}Ax\\
				&-x^\top A^\top S_{k+1}B (R+B^\top S_{k+1}B)^{-1} B^\top S_{k+1}Ax.
			\end{aligned}
			\]
			
			\medskip
			
			Поскольку это равенство должно выполняться для любого $x$,
			матрицы квадратичных форм должны быть равны.
			Получаем дискретную рекурсию Риккати:
			\[
			\boxed{
				\begin{aligned}
					S_k={}&Q+A^\top S_{k+1}A-A^\top S_{k+1}B
					(R+B^\top S_{k+1}B)^{-1} B^\top S_{k+1}A
				\end{aligned}
			}
			\]
			
			\medskip
			
			Терминальное условие следует из
			$V_N(x)=x^\top Q_Nx$:
			\[
			S_N=Q_N.
			\]
			
			Начиная с $S_N$, вычисляем $S_{N-1}$,
			затем $S_{N-2}$ и продолжаем в обратном направлении до $S_0$.
			
		\end{flushleft}
	\end{frame}
	%------------------------------------------------
	
	
	
	
	%------------------------------------------------
	\begin{frame}{Бесконечно-горизонтная стоимость}
		\begin{flushleft}
			
			Рассмотрим другую функцию стоимости:
			\[
			J=\sum_{k=0}^{\infty}
			\left(x_k^\top Qx_k+u_k^\top Ru_k\right).
			\]
			
			\medskip
			
			Заметим два важных отличия от конечно-горизонтного
			случая:
			\begin{itemize}
				\item Отсутствует терминальная стоимость.
				\item Стоимость накапливается за бесконечное
				число шагов по времени.
			\end{itemize}
			
			\medskip
			
			Это приводит к задаче \textbf{бесконечно-горизонтного
				дискретного LQR}.
			
		\end{flushleft}
	\end{frame}
	%------------------------------------------------
	
	\begin{frame}{Бесконечно-горизонтный случай: функция ценности}
		\begin{flushleft}
			
			В конечно-горизонтном случае функция ценности $V_k(x)$
			зависит от числа оставшихся шагов по времени.
			
			\medskip
			
			В бесконечно-горизонтном случае всегда остаётся бесконечно
			много шагов. Для стационарной системы и
			стоимости оптимальная стоимость до конца не зависит от
			текущего шага по времени.
			
			\medskip
			
			Поэтому можно записать
			\[
			V(x)=\min_{u_0,u_1,\ldots}
			\sum_{k=0}^{\infty}
			\left(x_k^\top Qx_k+u_k^\top Ru_k\right).
			\]
			
			\medskip
			
			Как и прежде, функция ценности является квадратичной:
			\[
			\boxed{V(x)=x^\top Sx.}
			\]
			
			Здесь $S$ — постоянная матрица, а не
			зависящая от времени последовательность матриц $S_k$.
			
		\end{flushleft}
	\end{frame}
	%------------------------------------------------
	
	\begin{frame}{Бесконечно-горизонтный дискретный LQR}
		\begin{flushleft}
			
			Уравнение Беллмана принимает вид
			\[
			V(x)=\min_u
			\left[
			x^\top Qx+u^\top Ru+V(Ax+Bu)
			\right].
			\]
			
			Подставляя $V(x)=x^\top Sx$, получаем
			\[
			x^\top Sx=\min_u
			\left[
			x^\top Qx+u^\top Ru
			+(Ax+Bu)^\top S(Ax+Bu)
			\right].
			\]
			
			\medskip
			
			Повторяя шаги из конечно-горизонтного случая, получаем
			оптимальный закон управления:
			\[
			\boxed{
				u=-Kx,
				\qquad
				K=(R+B^\top SB)^{-1}B^\top SA.
			}
			\]
			
			\medskip
			
			Подстановка оптимального управления в уравнение Беллмана
			даёт дискретное алгебраическое уравнение
			Риккати:
			\[
			\boxed{
				\begin{aligned}
					S= Q+A^\top SA -A^\top SB(R+B^\top SB)^{-1}B^\top SA.
				\end{aligned}
			}
			\]
			
		\end{flushleft}
	\end{frame}
	%------------------------------------------------
	
	\begin{frame}{Решение уравнения Риккати}
		\begin{flushleft}
			
			Дискретное алгебраическое уравнение Риккати можно
			решить численно с помощью стандартных библиотек теории управления.
			
			\medskip
			
			\begin{itemize}
				\item \textbf{Python (SciPy):}
				\texttt{scipy.linalg.solve\_discrete\_are(A, B, Q, R)}
				
				\medskip
				
				\item \textbf{MATLAB:}
				\texttt{dlqr(A, B, Q, R)}
				
				\medskip
				
				\item \textbf{Julia:}
				\texttt{lqr(Discrete, A, B, Q, R)}
			\end{itemize}
			
			\medskip
			
			После получения $S$ оптимальный коэффициент обратной связи равен
			\[
			K=(R+B^\top SB)^{-1}B^\top SA.
			\]
			
			Полученный закон управления $u=-Kx$ является стационарным.
			
		\end{flushleft}
	\end{frame}
	%------------------------------------------------

\begin{frame}{Литература}
	\begin{itemize}
		
		\item MIT OpenCourseWare:
		\bref{https://ocw.mit.edu/courses/16-323-principles-of-optimal-control-spring-2008/pages/lecture-notes/}
		{Principles of Optimal Control: Lecture 3, Dynamic Programming and Discrete LQR.}
		
		\item Russ Tedrake -- MIT:
		\bref{https://underactuated.mit.edu/lqr.html}
		{Underactuated Robotics: Linear Quadratic Regulators, including the discrete-time Riccati equation.}
		
		\item University of Washington:
		\bref{https://faculty.washington.edu/ratliffl/teaching/2020_LinearSystems.pdf}
		{Lecture 10: Linear Quadratic Regulator -- Discrete Time.}
		
		\item SciPy documentation:
		\bref{https://docs.scipy.org/doc/scipy/reference/generated/scipy.linalg.solve_discrete_are.html}
		{Linear Algebra: \texttt{solve\_discrete\_are}, numerical solution of the discrete-time algebraic Riccati equation.}
		
	\end{itemize}
\end{frame}





\myqrframe




\end{document}
